\section{Linear system for the core--shell particle}\label{app:sphere}
Each term of \cref{eq:fields} has velocity, pressure and traction amplitudes $(\mathcal U,\mathcal V,\mathcal P,\mathcal
S_{rr},\mathcal S_{r\theta})$ as defined in \cref{sec:sphere}. For the vortical terms they follow from
\cref{eq:vort_radial} with $g=\mathcal K_j$ or $\mathcal I_j$ and $\mathcal P=0$. For the potential terms, with the
pressure coefficients $c_j$ of \cref{lem:decomp}, they are
\begin{align*}
  \be:&\quad \mathcal U=1,\ \mathcal V=-1,\ \mathcal P=c_j\,r,\\
  \nabla(\be\cdot\nabla r^{-1}):&\quad \mathcal U=2r^{-3},\ \mathcal V=r^{-3},\ \mathcal P=-c_j\,r^{-2}.
\end{align*}
The rigid part $\be$ of \cref{eq:field_shell} carries the pressure $i\Wo^2\be\cdot\bx$, so its amplitudes at $r=1$ are
$\mathcal U=1$, $\mathcal V=-1$, $\mathcal S_{rr}=-i\Wo^2$ and $\mathcal S_{r\theta}=0$. Let a tilde denote the shell
field without its rigid part, and let shell tractions be computed with $\mu_s/\mu$. The conditions
\cref{eq:bc_core,eq:bc_shell} then read
\begin{align*}
  r=\kappa:&\quad \tilde{\mathcal U}^{(2)}=0,\qquad \tilde{\mathcal V}^{(2)}-\frac{\hat\ell_c}{\mu_s/\mu}\,\tilde{\mathcal S}^{(2)}_{r\theta}=0,\\
  r=1:&\quad \mathcal U^{(1)}=\mathcal U^{(2)},\quad \mathcal V^{(1)}=\mathcal V^{(2)},\quad
  \mathcal S^{(1)}_{rr}=\mathcal S^{(2)}_{rr},\quad
  \mathcal S^{(2)}_{r\theta}-\mathcal S^{(1)}_{r\theta}=\beta_s\alpha_s\left(\mathcal V^{(1)}+1\right).
\end{align*}
The kernels are rescaled to avoid overflow when $|\lambda_s|$ or $|\lambda|$ is large: $\mathcal K_s(r)e^{\lambda_s\kappa}$,
$\mathcal I_s(r)e^{-\lambda_s}$, $\mathcal K_f(r)e^{\lambda}$ and $\mathcal I_f(r)e^{-\lambda}$. The scaling factors are
removed from the coefficients after the solve. For a homogeneous porous particle ($\kappa=0$, $\beta_s=0$,
$\mu_s=\mu$, $\Wo\to0$) the solver reproduces the Debye--Bueche--Brinkman drag ratio
$\Omega=2\alpha_s^2(1-\alpha_s^{-1}\tanh\alpha_s)/[2\alpha_s^2+3(1-\alpha_s^{-1}\tanh\alpha_s)]$
(\cref{tab:validation}).

\section{Interface system in Fourier space}\label{app:interface}
For a plane mode $\propto e^{i\bk\cdot\bx_\parallel+sz}$ with velocity amplitude $\hat\bu$ and pressure amplitude
$\hat p$, the tractions at $z=0$ in a medium of viscosity $\mu_j$ are
\begin{equation*}
  \hat\sigma_{xz}=\mu_j\left(s\hat u_x+ik_x\hat u_z\right),\quad
  \hat\sigma_{yz}=\mu_j\left(s\hat u_y+ik_y\hat u_z\right),\quad
  \hat\sigma_{zz}=-\hat p+2\mu_js\hat u_z .
\end{equation*}
The potential mode has $\hat\bu=(ik_x,ik_y,s)$ and $\hat p=c_j$, with $s=-k$ (reflected) or $s=k$ (transmitted). The
poloidal mode has $\hat\bu=(ik_xs,ik_ys,k^2)$ and the toroidal mode $\hat\bu=(ik_y,-ik_x,0)$, both with $\hat p=0$ and
with $s=-q$ (reflected) or $s=q_p$ (transmitted). With superscripts ``inc'', $R$ and $T$ for the incident, reflected and
transmitted fields, the six rows of \cref{eq:interface_system} are
\begin{align*}
  &\hat\bu^{R}-\hat\bu^{T}=-\hat\bu^{\rm inc}\quad(3\ \text{rows}),\qquad
  \hat\sigma^{R}_{zz}-\hat\sigma^{T}_{zz}=-\hat\sigma^{\rm inc}_{zz},\\
  &\hat\sigma^{T}_{\tau z}-\hat\sigma^{R}_{\tau z}-\beta_p\alpha_p\hat u^{R}_\tau
   =\hat\sigma^{\rm inc}_{\tau z}+\beta_p\alpha_p\hat u^{\rm inc}_\tau\quad(\tau=x,y).
\end{align*}
The columns of $\mathsf A$ are normalised to unit Euclidean norm before solving, which removes the $O(k)$ degeneracy of
the mode vectors as $k\to0$. For a rigid wall the system reduces to $\hat\bu^R=-\hat\bu^{\rm inc}$.
